\documentclass[10pt,a4paper]{amsart}
\usepackage[margin=19mm]{geometry}
\usepackage{amsmath,amssymb,mathtools}
\usepackage[scr=rsfs,cal=euler]{mathalfa}
\usepackage{xfrac}
\usepackage[unicode,hidelinks,bookmarks=false]{hyperref}
\newcommand{\ie}{i.e.\ }
\newcommand{\cf}{cf.\ }
\newcommand{\resp}{respectively\ }
\newcommand{\Subsection}{\subsection}
\providecommand{\nobreakdash}{\nobreak\hbox{-}}
\setlength{\emergencystretch}{2em}
\multlinegap0pt
\usepackage{tikz}
\usetikzlibrary{cd}
\usepackage{amssymb}
\usepackage{tikz}
\newtheorem{lemma}{Lemma}
\newcommand{\rk}{\mathrm{rk}\,}
\newcommand{\set}[1]{\left\lbrace #1\right\rbrace}
\newcommand{\simto}{\xrightarrow{\raisebox{-0.7ex}[0ex][0ex]{$\sim$}}} % taken from https://tex.stackexchange.com/questions/479322/too-much-vertical-space-above-xrightarrow/479324#479324
\title{Interconnection of port-Hamiltonian systems is not dynamical}
\author{Jonas Kirchhoff}
\date{Source: arXiv:2608.29731v1 (CC BY 4.0)}
\begin{document}
\maketitle

\noindent\emph{Source and licence.} Interconnection of port-Hamiltonian systems is not dynamical. Version arXiv:2608.29731v1 (CC BY 4.0), distributed by arXiv under the Creative Commons Attribution 4.0 International licence, http://creativecommons.org/licenses/by/4.0/, as recorded in the arXiv licence metadata for this exact version and shown on the abstract page https://arxiv.org/abs/2608.29731v1. Full licence notice: \texttt{licenses/arxiv-2608.29731-CC-BY-4.0.txt}.\par\smallskip

\noindent\textit{Editorial scope.} Counted unit: the article's lemma lem:non\_constant\_case (source0.tex lines 1237--1257). The article's complete original proof of it is delivered with it (source0.tex lines 1258--1314, one \texttt{proof} environment of 57 lines). No mathematics is written, repaired or re-derived: the statement and the proof are the article's own lines, in the article's own notation, and no retained line is substituted or re-flowed.  Retained uncounted as the source-local context the proof needs: lines 838--840.  Nothing was omitted from any retained span, so the package carries no omission record.  No retained line contains a citation, so the wrapper carries no bibliography and no bibliography entry.  The retained lines contain the article's own diagram code, which the wrapper typesets with the standard tikz packages; no image file is delivered and the package declares no asset dependency.  Editorial formatting only, in addition: an AMS \texttt{amsart} preamble that restates the article's own declarations and macros used by the retained lines, namely the environments \texttt{lemma} on separate counters, together with the standard packages those lines need.  The retained environments print in this document as Lemma 1 in order of appearance, whereas the article prints the same items with the numbering of its own full document.  The complete native source of the article as distributed in the arXiv e-print for this exact version is bundled unchanged as \texttt{sources/208884/source0.tex}.
\begin{align}\label{eq:intersection}
\mathcal{D}_1\cap\mathcal{D}_2 = \set{(0_E(x),v,0_{E^*}(x),w)\in\mathcal{D}_1\,\big\vert\,(v,w)\in\mathcal{D}_2}
\end{align}

\begin{lemma}\label{lem:non_constant_case}
Let $\pi:E\to M$ and $F\to M$ be vector bundles over the same manifold, and let $\mathcal{D}_1\leq (E\oplus F)\oplus(E^*\oplus F^*)$ and $\mathcal{D}_2\leq F\oplus F^*$ be Dirac structures. If $\mathcal{D}_1\cap\mathcal{D}_2$ defined in~\eqref{eq:intersection} has constant rank, then $\mathcal{C}^k([t_0,t_1],\mathcal{D}_1\circ\mathcal{D}_2)$ and
\begin{align*}
\set{(X,\alpha)\in\mathcal{C}^k([a,b],E\oplus E^*)\,\left\vert
\begin{array}{l}
\exists (Y,\beta)\in\mathcal{C}^k([t_0,t_1],D_2):\\
\hspace*{1.5cm}(X,Y,\alpha,\beta)\in\mathcal{C}^k([t_0,t_1],D_1)
\end{array}
\right.}
\end{align*}
coincide for all $t_0<t_1\in\mathbb{R}$ and $k\in\mathbb{N}$; likewise, $\mathcal{L}^p_c([t_0,t_1],,\mathcal{D}_1\circ\mathcal{D}_2)$ and
\begin{align*}
\set{(X,\alpha)\in\mathcal{L}^p_c([a,b],E\oplus E^*)\,\left\vert
\begin{array}{l}
\exists (Y,\beta)\in\mathcal{L}^p_c([t_0,t_1],D_2):\\
\hspace*{1.5cm}(X,Y,\alpha,\beta)\in\mathcal{L}^p_c([t_0,t_1],D_1)
\end{array}
\right.}
\end{align*}
coincide for all $p\in[1,\infty]$.
\end{lemma}

\begin{proof}
Put $G:=p_{F\oplus F^*}\left(\mathcal{D}_1\cap\mathcal{D}_2\right)$, and let $(X,\alpha)\in \mathcal{C}^k([t_0,t_1],\mathcal{D}_1\circ\mathcal{D}_2)$ for some $t_0<t_1\in\mathbb{R}$ and $k\in\mathbb{N}$. There exists $\varepsilon\in (0,\infty)$ and representatives (or extensions) $(\overline{X},\overline{\alpha})\in\mathcal{C}^k((t_0-\varepsilon,t_1+\varepsilon),\mathcal{D}_1\circ\mathcal{D}_2)$. Put $x := \pi_E\circ \overline{X}$. Since $(t_0-\varepsilon,t_1+\varepsilon)$ is contractible, we find trivialisations
\begin{align*}
x^*(E\oplus E^*) & \simto (t_0-\varepsilon,t_1+\varepsilon)\times\mathbb{R}^{2n},\\
\varphi:x^*\mathcal{D}_1 & \simto (t_0-\varepsilon,t_1+\varepsilon)\times\mathbb{R}^{n+m},
\end{align*}
and trivialisations of $x^*G$, $x^*\mathcal{D}_2$ and $x^*(F\oplus F^*)$ so that the diagram
\begin{equation}\label{eq:comm_diagram}
\begin{tikzcd}
x^*G \arrow[d,"\subseteq"]\arrow[r,"\sim"] & (t_0-\varepsilon,t_1+\varepsilon)\times\mathbb{R}^\ell\arrow[d,"\subseteq"]\\
x^*\mathcal{D}_2\arrow[r,"\psi"]\arrow[d,"\subseteq"] & (t_0-\varepsilon,t_1+\varepsilon)\times\mathbb{R}^\ell\oplus\mathbb{R}^{m-\ell}\arrow[d,"\subseteq"]\\ 
x^*(F\oplus F^*)\arrow[r,"\sim"] & (t_0-\varepsilon,t_1+\varepsilon)\times\mathbb{R}^\ell\oplus\mathbb{R}^{m-\ell}\oplus\mathbb{R}^{m}\\ 
\end{tikzcd}
\end{equation}
commutes. In particular, we have a kernel representation
\begin{align*}
x^*\mathcal{D}_1 \simeq \set{(t,X,Y,\alpha,\beta)\in\mathbb{R}^{2(n+m)}\,\left\vert\, A(t)\begin{pmatrix}
X\\\alpha
\end{pmatrix} = B(t)\begin{pmatrix}
Y\\\beta
\end{pmatrix}\right.}
\end{align*}
for matrix fields
\begin{align*}
A & \in\mathcal{C}^\infty\big((t_0-\varepsilon,t_1+\varepsilon),\mathbb{R}^{(n+m)\times 2n}\big),\\
B = [B_1, B_2, B_3] & \in\mathcal{C}^\infty\big((t_0-\varepsilon,t_1+\varepsilon),\mathbb{R}^{(n+m)\times (\ell+(m-\ell)+m)}\big).
\end{align*}
Therefore, we have
\begin{align*}
x^*(\mathcal{D}_1\circ\mathcal{D}_2) \simeq \set{(t,X,\alpha)\in\mathbb{R}^{2n}\,\left\vert\,\begin{array}{l}\exists z_1\in\mathbb{R}^{\ell}~\exists z_2\in\mathbb{R}^{m-\ell}:\\
\hspace*{1cm} A(t)\begin{pmatrix}
X\\\alpha
\end{pmatrix} = \begin{bmatrix}B_1(t), B_2(t)\end{bmatrix}\begin{pmatrix}
z_1\\z_2
\end{pmatrix}
\end{array}\right.}
\end{align*}
and we observe that
\begin{align*}
& \ker \begin{bmatrix}B_1(t), B_2(t)\end{bmatrix}\\
&\quad \simeq \set{(Z,\beta)\in(\mathcal{D}_2)_{x(t)}\,\big\vert\,\forall (X,\alpha)\in(\mathcal{D}_1\circ\mathcal{D}_2)_{x(t)}:(X,Y,\alpha,\beta)\in\mathcal{D}_1} = G_{x(t)}
\end{align*}
for all $t\in(t_0-\varepsilon,t_1+\varepsilon)$ and thus the commutativity of~\eqref{eq:comm_diagram} yields $B_2\equiv 0$ and $\rk B_1\equiv \ell$. Hence,  $B_1$ admits the smooth left-inverse
\begin{align*}
M := (B_1(\cdot)^\top B_1(\cdot))^{-1}B_1(\cdot)^\top\in\mathcal{C}^\infty\big((t_0-\varepsilon,t_1+\varepsilon),\mathbb{R}^{\ell\times(n+m)}\big).
\end{align*}
Define
\begin{align*}
(\overline{Y},\overline{\beta}):(t_0-\varepsilon,t_1+\varepsilon) & \to F\oplus F^*,\\
t & \mapsto \psi^{-1}\left(M(t)A(t)\varphi\left(\overline{X}(t),\overline{\alpha}(t)\right)\right)
\end{align*}
which is in $\mathcal{C}^k\big((t_0-\varepsilon,t_1+\varepsilon),\mathcal{D}_2\big)$ and satisfies
\begin{align*}
\forall t\in (t_0-\varepsilon,t_1+\varepsilon): \left(\overline{X}(t),\overline{Y}(t),\overline{\alpha}(t),\overline{\beta}(t)\right)\in\mathcal{D}_1.
\end{align*}
Taking the restriction on $[t_0,t_1]$ gives the claimed identity. The second part of the lemma is proven analogously.
\end{proof}

\end{document}
